%! Inverse Functions — Unit 5, Lesson 5.7 — Lesson Plan
\documentclass[10pt]{article}
\usepackage{algebra2-article}
\usepackage{algebra2-boxes}
\usepackage{graphicx}   % -article does NOT load graphicx; needed for thumbnails/screenshots
\graphicspath{{images/}}
\newcommand{\TallMath}[1]{$\displaystyle #1\rule[-1.4em]{0pt}{3.2em}$}
\newcommand{\CourseName}{Algebra 2}
\newcommand{\MeetingLength}{55 minutes}
\newcommand{\UnitNumberName}{Unit 5: Radical Functions \quad}
\newcommand{\LessonNumberName}{Lesson 5.7: Inverse Functions}

\begin{document}
\begin{center}
    {\Large\bfseries \CourseName \\
    {\small\bfseries \UnitNumberName \LessonNumberName}}
\end{center}

\begin{tcolorbox}[colback=forestbg, colframe=forest, boxrule=0.9pt, arc=2mm,
    left=2mm, right=2mm, top=1.5mm, bottom=1.5mm]
    \textbf{Primary Objective:} Students can \emph{justify} that two functions are inverses by
    composing them in \textbf{both} orders and obtaining $f\big(g(x)\big)=g\big(f(x)\big)=x$, and can
    \emph{build} the inverse of a linear, cubic, quadratic, or square root function by
    \textbf{swapping and solving} --- reversing the operations \emph{and the order}. They can read
    points of $f^{-1}$ off a table or a graph of $f$, explain the graph of $f^{-1}$ as the
    \textbf{reflection over $y=x$} with domain and range trading places, apply the
    \textbf{horizontal line test} to decide whether an inverse is a function, and state the
    \textbf{restriction} an inverse of a radical or quadratic requires --- naming the \emph{range of
    $f$} as its source. This is the \textbf{unit capstone}: it closes the question opened in Lesson
    5.0 by proving that $y=x^{2}$ has two arms and $y=x^{3}$ has one, from which every contrast
    between the square root and cube root families follows.
    \\[2pt]
    \textbf{Standards (2023 VA SOL):} \textbf{A2.F.2i} (governing --- determine the inverse of a
    function algebraically and graphically for linear, quadratic, and square root functions;
    \emph{justify and explain why two functions are inverses of each other}, which is the
    composition test in Notes \S2 and is assessed in every component) and \textbf{A2.F.2j} (graph
    the inverse as a reflection over $y=x$ --- Notes \S4, Activity Tier~R Part~3 and Tier~A Part~3,
    Homework item~4). It rests directly on \textbf{A2.F.2k} (Lesson 5.6's composition, now used as
    a proof tool) and revisits \textbf{A2.F.2a} (domain and range, traded), \textbf{A2.EO.2a/c}
    (roots undoing powers; $\sqrt{x^{2}}=|x|$), \textbf{A2.EI.5a} (the solve-for-the-other-variable
    move), and \textbf{A2.F.2f} (interpret in context --- the pendulum model).
\end{tcolorbox}

\renewcommand{\arraystretch}{1.4}
\begin{skillbox}[Priority Ideas \& Skills]{goldbox}
\begin{tabularx}{\linewidth}{@{}X|X@{}}
\textbf{Priority Ideas \& Skills} & \textbf{Key Understandings (the ``why'')} \\[2pt]
\hline
\vspace{-6pt}
\begin{itemize}[leftmargin=1.1em, nosep, after=\vspace{2pt}]
  \item Read $f^{-1}(x)$ correctly --- a label, never an exponent.
  \item \textbf{Justify} an inverse pair by composing \emph{both} ways and getting $x$ each time.
  \item Build $f^{-1}$ by \textbf{swap and solve}, for linear, cubic, quadratic, and square root functions.
  \item Reverse the operations \emph{and the order} --- and know why that is not optional.
  \item Find points of $f^{-1}$ from a table or a graph; recognize the \textbf{reflection over $y=x$}.
  \item Swap domain and range, and state the \textbf{restriction} an inverse needs.
  \item Apply the \textbf{horizontal line test} to decide whether an inverse function exists.
  \item Explain why $y=x^{2}$ needs a restricted domain and $y=x^{3}$ does not.
\end{itemize}
&
\vspace{-6pt}
\begin{itemize}[leftmargin=1.1em, nosep, after=\vspace{2pt}]
  \item ``Inverse'' is not new information --- students said it in their own words in 5.6 Tier A. Today it gets a name, a test, and a construction.
  \item \emph{Composition is the proof.} Yesterday's tool is today's evidence, which is exactly why 5.6 had to come first.
  \item \emph{One order is never enough.} $\left(\sqrt{x}\,\right)^{2}=x$ but $\sqrt{x^{2}}=|x|$ --- the same counterexample for the third day running.
  \item Shoes come off before socks. Rosa's $\frac59f-32$ is Warm-Up item 3 in disguise, and it is the exit ticket's most-chosen distractor.
  \item Swapping the equation and swapping every point are the same move --- which is why the picture is a reflection.
  \item \emph{The restriction is read off the range of $f$; it is never guessed.} An inverse with no restriction is a different function, not a sloppy one.
  \item Swap-and-solve will hand you a formula for a function that has no inverse. The horizontal line test is what licenses the word ``the.''
  \item \textbf{The capstone sentence:} $x^{2}$ has two arms, $x^{3}$ has one. Every difference between $\sqrt{x}$ and $\sqrt[3]{x}$ comes from that.
\end{itemize}
\end{tabularx}
\end{skillbox}

\begin{skillbox}[{Vocabulary, Concepts, \& Theorems}]{sky}
\renewcommand{\arraystretch}{1.3}
\begin{tabularx}{\linewidth}{@{}l X@{}}
\textbf{Inverse function} & A function that undoes another: if $f$ sends $a$ to $b$, then $f^{-1}$ sends $b$ back to $a$. \\
\textbf{$f^{-1}(x)$} & ``$f$ inverse of $x$.'' The $-1$ is a \emph{label}, not an exponent --- $f^{-1}(x)$ is never $1/f(x)$. \\
\textbf{The inverse test} & $f$ and $g$ are inverses exactly when $f(g(x))=x$ \emph{and} $g(f(x))=x$. Both orders are required. \\
\textbf{Swap and solve} & The construction: write $y=f(x)$, trade $x$ and $y$, solve for $y$, rename it $f^{-1}(x)$. \\
\textbf{Reverse the order} & To undo a sequence of steps you undo them backwards --- shoes before socks. The lesson's most under-taught idea. \\
\textbf{Reflection over $y=x$} & What swapping coordinates does to a graph: $(a,b)\mapsto(b,a)$. Points already on $y=x$ never move. \\
\textbf{Horizontal line test} & A graph has an inverse \emph{function} exactly when no horizontal line meets it more than once. \\
\textbf{Restricted domain} & Deliberately shrinking a domain (as with $y=x^{2}$, $x\ge0$) so a function stops repeating outputs and earns an inverse. \\
\textbf{Domain/range swap} & domain of $f^{-1}$ $=$ range of $f$, and range of $f^{-1}$ $=$ domain of $f$. The source of every restriction today. \\
\textbf{Involution} & (Tier E) A function that is its own inverse --- $6-x$, $\frac1x$ --- whose graph is symmetric about $y=x$. \\
\end{tabularx}
\end{skillbox}

\begin{fixedskillbox}[Activate Prior Knowledge \& Spiral Review]{forestbg}
    Three items, each a piece of today's lesson. \textbf{(1)} \emph{Solve for the other variable}
    --- $y=3x-7$, $y=\frac{x+5}{2}$, $y=x^{3}+1$. This is the engine of swap-and-solve, and students
    can already do it; part (c) is the unit's own contribution, since a cube root is what undoes a
    cube. \textbf{(2)} Lesson 5.6's Tier~A discovery, promoted to a warm-up: $f(x)=2x+6$ and
    $g(x)=\frac{x-6}{2}$ compose to $x$ in \emph{both} orders. \textbf{Collect the wordings students
    invented yesterday} (``they cancel,'' ``it goes backwards'') and put two or three on the board
    --- when \textbf{inverse} appears in \S1, the definition should look like a formal version of
    what the class already said. \textbf{(3)} ``Double, then add $5$, and I get $17$'' --- undo it,
    and name the order you undid in. Nobody halves first. Point back at that answer during the Hook,
    when half the room writes $\frac59f-32$; it is the same mistake somewhere much harder to see.
    Resolve none of the three.
\end{fixedskillbox}

\begin{skillbox}[Hook]{forestbg}
    Your app reads $^\circ$F, your cousin's reads $^\circ$C, and $F(c)=\frac95c+32$ goes one way.
    Fill the table on the board ($0\to32$, $100\to212$, $20\to68$), then put up two students' attempts
    at the \emph{other} direction: \textbf{Rosa} writes $C(f)=\frac59f-32$; \textbf{Malik} writes
    $C(f)=\frac59(f-32)$. Do not adjudicate --- \emph{test at $f=212$, where the room already knows
    the answer must be $100$}. Rosa gives $\approx85.8$; Malik gives $100$. Then ask the question
    that carries the lesson: \textbf{Rosa reversed both operations correctly. What did she not
    reverse?} She undid them in the same order they were done. Send the class straight back to
    Warm-Up item 3 --- nobody undoes ``double, then add $5$'' by halving first --- and name the rule
    you want repeated all period: \emph{to undo a process you reverse the operations \textbf{and the
    order}. Shoes before socks.} Then state the day's two questions: how do you build the other-way
    rule for \emph{any} function, and how do you \emph{prove} you got it right?
\end{skillbox}

\begin{skillbox}[Lesson]{forestbg}
\begin{multicols}{2}
\textbf{1. The word and the notation.} Put the machine diagram up --- $x\to f\to f^{-1}\to x$ ---
beside the students' own wordings from yesterday. Then spend a full minute on the red box, because
it is the only place all year this gets said: \textbf{the $-1$ is not an exponent}. Make it concrete
with $f(x)=2x$: the inverse is $\frac{x}{2}$, the reciprocal is $\frac{1}{2x}$, and at $x=4$ those
are $2$ and $\frac18$. Close by restating Warm-Up item 3 in the new language.

\textbf{2. The test comes before the construction.} This ordering is deliberate: students need a
verdict they trust before they start building partners of their own. $f(g(x))=g(f(x))=x$, both
orders, every time. Run (a) the Warm-Up pair, which passes; then (b) $h(x)=\frac{x}{2}-6$ against
$f(x)=2x+6$ --- right operations, wrong order, and the composition comes out $x-6$. Stress that it
is \emph{close}: an error off by a constant survives an eyeball check and only composition catches
it. Then (c), yesterday's cliffhanger: $\left(\sqrt{x}\,\right)^{2}=x$ but $\sqrt{x^{2}}=|x|$ ---
\emph{one order passing is not a verdict}, which is why the test names both.

\textbf{3. Swap and solve.} Four steps, and step 3 is just Warm-Up item 1. Work the table: $3x-7$,
$\frac{x}{4}+2$, $x^{3}+1$, $\sqrt{x-2}$. Rows (c) and (d) are the unit's own moves --- a cube root
undoes a cube, squaring undoes a square root --- so say the phrase \emph{inverse families} and point
back at 5.0. Row (d) is where the lesson gets hard: $f^{-1}(x)=x^{2}+2$ \textbf{for $x\ge0$}, and
that condition is not bookkeeping. Check it by asking what $f^{-1}(-3)$ would claim, when $f$ never
sent anything to $-3$.

\columnbreak

\textbf{4. The picture.} Swapping the equation and swapping every point are the same move: $(a,b)$
becomes $(b,a)$. Do the $\sqrt{x}$ table first --- the rows literally trade places, and it is the
table students built in 5.0 --- then the two graphs with $y=x$ dashed. Read the linear picture
aloud: $x$-intercept becomes $y$-intercept, and the graphs cross at $(4,4)$ \emph{on the mirror
line}, which is the only place a graph and its reflection can meet. That fact pays off in Tier E.

\textbf{5. The capstone --- protect ten minutes for this.} Three pre-drawn graphs answer the
question the unit opened with. $y=x^{2}$ fails the horizontal line test: $(-2,4)$ and $(2,4)$ swap
into $(4,-2)$ and $(4,2)$, one input with two outputs, which breaks the \emph{vertical} line test.
Restrict to $x\ge0$ and one arm survives --- and its inverse is $\sqrt{x}$. Meanwhile $y=x^{3}$
passes untouched, so $\sqrt[3]{x}$ needs no restriction. Land the sentence and leave it on the
board: \emph{$x^{2}$ has two arms, $x^{3}$ has one} --- and every difference between the two radical
families in this unit comes from that.

\textbf{6. Domain and range trade places.} Five minutes; fold into \S5 if short. The table for
$\sqrt{x-2}$ and $x^{2}+2$ shows the same two sets swapped, which retroactively explains where the
$x\ge0$ in \S3 came from: it is the \emph{range of $f$}, written down.
\end{multicols}
\end{skillbox}

\begin{skillbox}[Active Monitoring]{redbox}
Circulate during the notes practice and the tiers. Watch for and correct:
\begin{itemize}[leftmargin=1.2em, nosep]
  \item \textbf{$f^{-1}$ read as an exponent} --- $\frac{1}{f(x)}$. Rare but total: the student thinks the lesson is about reciprocals. Fix it one-on-one with $f(x)=2x$ evaluated at $4$. It is exit-ticket distractor C and Homework 5 (Priya);
  \item \textbf{undoing in the wrong order} --- the day's signature error. Rosa's $\frac59f-32$, $\frac{x}{2}-6$ in \S2(b), Activity Tier~R 2(c), Homework 1(c) and 5 (Devon), exit-ticket distractor B. Correct it with the sentence, not the algebra: \emph{which step happened last?}
  \item \textbf{an inverse announced without being tested.} Treat it as unfinished work. The whole point of \S2 preceding \S3 is that students own a verdict procedure;
  \item \textbf{the missing restriction} on the inverse of a radical --- $x^{2}+2$ with no $x\ge0$. \textbf{This is the one that costs test points.} It is not a detail: it is a different function that fails the test on every negative input;
  \item \textbf{guessing the restriction as $x\ge0$ by reflex.} Activity Tier~A row (f) needs $x\ge-3$, and a student who writes $x\ge0$ has pattern-matched. The question to ask is always \emph{what is the range of $f$?}
  \item \textbf{swap-and-solve applied to a function with no inverse} --- e.g.\ writing ``the inverse of $x^{2}-4$ is $\sqrt{x+4}$'' with no restriction. The algebra is fine; the sentence is false. Activity Tier~A J3;
  \item \textbf{sign slips when the variable is subtracted} --- $f(x)=7-2x$ (Homework 2(c)), where students reach for ``divide by $-2$'' first;
  \item \textbf{$\sqrt{x^{2}}$ written as $x$.} Third day running (6.1, 6.6, today), and today it is load-bearing: it is the reason the test needs both orders.
\end{itemize}
Cold-call prompts: ``Which step happened \emph{last} --- so which one comes off \emph{first}?''
``You built a partner. What is your evidence?'' ``Your inverse has $x\ge0$ on it. Where did that
number come from?'' ``$(2,2)$ didn't move when we reflected. Why not, and which other points would
do the same?'' ``This one simplified to a nice formula and still has a restriction. What is being
refused, and by which function?''
\end{skillbox}

\begin{skillbox}[Group Work \& Differentiation]{redbox}
\textbf{Undo It} activity (tiered; mirrors the notes).
\begin{multicols}{3}
\textbf{Tier R --- Remediate.} Part 1 is arithmetic-first and deliberately avoids the word
``inverse'' until the last column: list what $2x-1$ and $\frac{x}{5}+4$ do \emph{in order}, then undo
\emph{in reverse order}, then write the rule. Part 2 tests three pairs by composition, and row (c)
($3x+1$ against $\frac{x}{3}+1$) is the whole lesson in one line --- right operations, wrong order,
caught only by composing. Part 3 swaps a table and reads four marked points off a pre-drawn graph;
$(2,2)$ does not move, which is the observation Tier A J4 and Tier E Part 1 both need.

\columnbreak
\textbf{Tier A --- Approaching.} Part 1 is six swap-and-solves, two of which need a restriction ---
and \textbf{row (f), $\sqrt{x}-3$, is the assessed discriminator}, since its condition is $x\ge-3$
and reflex says $x\ge0$. Part 2 is the A2.F.2i justification, both orders written out. Part 3 runs
the horizontal line test on four pre-drawn graphs (line, parabola, square root, and a Unit~3 cubic
with two turning points, where the worst line crosses three times). Part 4: \textbf{J1} asks where
the $-3$ comes from (the range of $f$ --- circulate for this sentence); \textbf{J3} is the deepest
item on the page, since swap-and-solve will happily produce a formula for a function that has no
inverse.

\columnbreak
\textbf{Tier E --- Extension.} Four parts. \emph{(1)} Finish the Hook: build $F^{-1}$, prove it both
ways, then solve $F(c)=c$ to find $-40^{\circ}$ --- the one temperature that reads the same on both
scales --- and explain it with the mirror line rather than algebra. \emph{(2)} Functions that are
their own inverse: $6-x$, $\frac1x$, and $\frac{2x+3}{x-2}$ (a Unit~4 callback), whose graphs are
symmetric about $y=x$. \emph{(3)} Socks and shoes in symbols: $(f\circ g)^{-1}=g^{-1}\circ f^{-1}$,
discovered by testing both orders against $2x+6$. \emph{(4)} Restrict $y=x^{2}$ to $x\le0$ instead
and get $-\sqrt{x}$, so ``\emph{the} inverse'' is not well posed until the restriction is named ---
closing on $A=\pi r^{2}$, where the situation chooses the branch for you.
\end{multicols}
\end{skillbox}

\begin{skillbox}[Individual Work \& Assessment]{redbox}
\textbf{Exit Ticket} (independent, no notes): verify $f(x)=4x-8$ against $g(x)=\frac{x+8}{4}$ in both
orders (each gives $x$), with 1(d) --- \emph{why does the test require both orders?} --- as the
conceptual item, whose complete answer produces the counterexample rather than restating the rule;
then build $f^{-1}$ for $f(x)=\sqrt{x-1}$, obtaining $x^{2}+1$ \textbf{with $x\ge0$}, and justify
the restriction by naming the \emph{range of $f$}; then an \textbf{SOL-style multiple-choice} item,
$f(x)=3x+12$, with options $\frac{x-12}{3}$, $\frac{x}{3}-12$, $\frac{1}{3x+12}$, and $3x-12$
(answer A). The three distractors are the lesson's three named errors, one each: \textbf{B} reversed
the operations but not the order (Rosa), \textbf{C} read the $-1$ as an exponent, \textbf{D} changed
a sign and undid nothing. Collect and sort into ``got it / order reversed / read as an exponent /
domain dropped.'' \textbf{This is the last exit ticket before the unit test, so the piles are also
the review-day agenda} --- and the fourth pile is the expensive one, since a bare formula earns no
A2.F.2i justification credit.
\end{skillbox}

\begin{skillbox}[Reinforcement \& Extension]{goldbox}
\textbf{Homework:} four verification pairs with one impostor (row (c) has the right operations in the
wrong order, and (e) asks students to name \emph{which} of the two things went wrong); six
swap-and-solves, of which (f) needs $x\ge4$ and (c), $7-2x$, is the sleeper sign trap; a table item
ending with the reflection rule applied to a point not in the table; a two-graph item combining the
reflection reads with a horizontal line test verdict and a repair; an error analysis carrying the
day's two \emph{distinct} misconceptions side by side --- Priya takes a reciprocal, Devon reverses
the operations but not the order --- each caught with a single number; and a two-part justification
item that is the capstone of the unit. \textbf{Extension:} the pendulum, $T(L)=2\pi\sqrt{L/32}$,
inverted to $L(T)=\frac{8T^{2}}{\pi^{2}}$. It is the only place in the unit where the \emph{inverse}
is what a real person actually needs --- a clockmaker never asks how fast a given pendulum swings,
they ask how long to make it --- and a $2$-second swing gives $\approx3.24$ feet, i.e.\ about $39$
inches, which students can check against any real grandfather clock. Part (b) rediscovers that
quadrupling the length doubles the period ($\sqrt{4L}=2\sqrt{L}$, the flattening first read off the
graph in 5.0), and part (f) is the quiet one: every restricted domain this week had to be argued
for, and here the situation supplies it for free.
\textbf{Preview:} the \textbf{Unit 5 test}. This lesson closes the unit --- rational exponents,
simplifying and operating, rationalizing, graphing and transforming, solving with extraneous checks,
composition, and now inverses. The homework's closing note ties the whole arc back to the question
Lesson 5.0 opened: $x^{2}$ has two arms and $x^{3}$ has one, and every difference between $\sqrt{x}$
and $\sqrt[3]{x}$ follows from that one sentence. Assign the practice test with the review day.
\end{skillbox}
\end{document}
